\documentclass[11pt]{article}
\usepackage{amsmath,amssymb}

\setlength{\parindent}{0pt}
\setlength{\parskip}{0.7em}
\setlength{\textwidth}{6.5in}
\setlength{\oddsidemargin}{0in}
\setlength{\evensidemargin}{0in}
\setlength{\topmargin}{-0.4in}
\setlength{\textheight}{9.2in}

\begin{document}

\begin{center}
{\Large SYSC 5108 Exam Study: Assignment 4, Question 2}\\
\vspace{0.3em}
{\large Markov chain sender example: marginals, revisits, and time averages}\\
\vspace{0.3em}
Source question: \texttt{SYSC 5108 W26 Assignment 4.pdf}, Question 2
\end{center}

\section*{Question Summary}

Question 2 refers to the example on Slide 10 of Chapter 6. The sender has a two-state Markov chain with transition matrix
\[
P=
\begin{bmatrix}
p & 1-p\\
0 & 1
\end{bmatrix}.
\]
Interpretation:
\begin{itemize}
\item State \(0\): the sender is still trying to send the message.
\item State \(1\): the message has been sent successfully.
\end{itemize}

The chain starts in state \(0\) at time \(n=0\), so the initial distribution is
\[
\pi(0)=
\begin{bmatrix}
1 & 0
\end{bmatrix}.
\]

State \(1\) is absorbing, because once the message is sent successfully, the sender stays in state \(1\):
\[
P_{11}=1.
\]

\section*{Step 1: Use the Marginal Distribution Update Rule}

Chapter 6 slides define the marginal distribution by
\[
\pi_i(n)=P\{X_n=i\}
\]
and the update rule
\[
\pi^T(n+1)=\pi^T(n)P.
\]

The textbook also describes a Markov chain through a transition operator and explains that repeated application of the transition matrix updates the state distribution over time.

Let
\[
\pi(n)=
\begin{bmatrix}
\pi_0(n) & \pi_1(n)
\end{bmatrix}.
\]
Then
\[
\pi(n+1)=\pi(n)
\begin{bmatrix}
p & 1-p\\
0 & 1
\end{bmatrix}.
\]

\section*{Part (a): Find \(\pi_0(n)\) and \(\pi_1(n)\)}

From the update rule:
\[
\pi_0(n+1)=p\,\pi_0(n),
\]
because the only way to be in state \(0\) at time \(n+1\) is to already be in state \(0\) at time \(n\) and remain there.

Since
\[
\pi_0(0)=1,
\]
we get
\[
\pi_0(1)=p,\qquad
\pi_0(2)=p^2,\qquad
\ldots
\]
So by induction,
\[
\boxed{\pi_0(n)=p^n.}
\]

Because the probabilities must sum to 1,
\[
\pi_1(n)=1-\pi_0(n)=1-p^n.
\]
So
\[
\boxed{\pi_1(n)=1-p^n.}
\]

\textbf{Check:}
\[
\pi(1)=
\begin{bmatrix}
1 & 0
\end{bmatrix}
\begin{bmatrix}
p & 1-p\\
0 & 1
\end{bmatrix}
=
\begin{bmatrix}
p & 1-p
\end{bmatrix},
\]
which matches the formula for \(n=1\).

\section*{Part (b): Probability the Sender Will Revisit State 0}

Because state \(1\) is absorbing, the chain can revisit state \(0\) only if it stays in state \(0\) on the first step.

Starting from state \(0\):
\[
P(X_1=0\mid X_0=0)=p.
\]
If the chain moves to state \(1\), it never returns to state \(0\). Therefore,
\[
\boxed{P(\text{sender revisits state }0)=p.}
\]

\textbf{Important distinction:}
\begin{itemize}
\item Probability of at least one revisit: \(p\).
\item Expected number of revisits (from the regenerative method on the slide): \(\eta(0)=\frac{p}{1-p}\) for \(0\le p<1\).
\end{itemize}
The question asks for the first quantity, not the second.

\section*{Part (c): Time Average Probabilities \(\bar\pi_0\) and \(\bar\pi_1\)}

The time-average probability of state \(i\) is
\[
\bar\pi_i
=
\lim_{N\to\infty}\frac{1}{N}\sum_{n=0}^{N-1}\pi_i(n).
\]

Using part (a):
\[
\bar\pi_0
=
\lim_{N\to\infty}\frac{1}{N}\sum_{n=0}^{N-1}p^n.
\]
For \(0\le p<1\), the geometric sum is bounded:
\[
\sum_{n=0}^{N-1}p^n=\frac{1-p^N}{1-p}.
\]
Therefore
\[
\bar\pi_0
=
\lim_{N\to\infty}
\frac{1-p^N}{N(1-p)}
=0.
\]
So
\[
\boxed{\bar\pi_0=0.}
\]

Since
\[
\pi_1(n)=1-p^n,
\]
we get
\[
\bar\pi_1
=
\lim_{N\to\infty}\frac{1}{N}\sum_{n=0}^{N-1}(1-p^n).
\]
This is
\[
\bar\pi_1
=
\lim_{N\to\infty}
\left(
1-\frac{1}{N}\sum_{n=0}^{N-1}p^n
\right)
=1-0=1.
\]
So
\[
\boxed{\bar\pi_1=1.}
\]

\textbf{Interpretation:} for \(0\le p<1\), the sender spends only a transient amount of time in state \(0\), and eventually stays forever in the absorbing success state \(1\).

\textbf{Edge case:} if \(p=1\), then the message is never sent and the chain stays in state \(0\) forever, so
\[
\bar\pi_0=1,\qquad \bar\pi_1=0.
\]
In the usual episodic sender example, we assume \(0\le p<1\).

\section*{Part (d): Probability the Sender Is in State 0 Before Success}

Before the message is sent successfully, the chain has not yet entered state \(1\). Since the only states are \(0\) and \(1\), the sender must be in state \(0\).

So, for any time before successful transmission,
\[
\boxed{P(\text{sender is in state }0\text{ before success})=1.}
\]

Another way to state this is:
\[
X_n=0 \quad \text{for all } n<T,
\]
where \(T\) is the first time the message is sent successfully.

\section*{Step 2: Why the Results Make Sense}

This chain has one transient state and one absorbing state:
\begin{itemize}
\item State \(0\) is transient when \(0\le p<1\), because eventually the chain leaves it.
\item State \(1\) is absorbing and therefore recurrent.
\end{itemize}

That explains all four answers:
\begin{itemize}
\item \(\pi_0(n)=p^n\) decays geometrically.
\item \(\pi_1(n)=1-p^n\) increases to 1.
\item the probability of ever revisiting state \(0\) is just the probability of staying there on the next step, namely \(p\),
\item the time average concentrates entirely on state \(1\).
\end{itemize}

\section*{Concise Exam Answer}

With
\[
P=
\begin{bmatrix}
p & 1-p\\
0 & 1
\end{bmatrix},
\qquad
\pi(0)=
\begin{bmatrix}
1 & 0
\end{bmatrix},
\]
the marginals satisfy
\[
\pi_0(n+1)=p\,\pi_0(n),\qquad \pi_0(0)=1,
\]
so
\[
\boxed{\pi_0(n)=p^n,\qquad \pi_1(n)=1-p^n.}
\]
The sender revisits state \(0\) only if it remains in state \(0\) on the first step, so
\[
\boxed{P(\text{revisit state }0)=p.}
\]
For \(0\le p<1\),
\[
\bar\pi_0=\lim_{N\to\infty}\frac{1}{N}\sum_{n=0}^{N-1}p^n=0,
\qquad
\bar\pi_1=1.
\]
Hence
\[
\boxed{\bar\pi_0=0,\qquad \bar\pi_1=1.}
\]
Before the message is sent successfully, the sender must still be in state \(0\), so
\[
\boxed{P(\text{sender is in state }0\text{ before success})=1.}
\]

\section*{Cross References Used}

\begin{itemize}
\item Assignment: \texttt{SYSC 5108 W26 Assignment 4.pdf}, Question 2.
\item Local submitted Assignment 4 PDF checked for comparison, Question 2 response.
\item Slides: Chapter 6 slide deck, slides 6--10 on Markov chains, marginal distributions, stationary transitions, and the episodic sender example.
\item Slides: Chapter 6 slide deck, slide 11 on classification of states.
\item Textbook: Goodfellow, Bengio, and Courville, \textit{Deep Learning}, Chapter 17 on Markov chains, transition operators, stochastic matrices, and stationary distributions.
\end{itemize}

\end{document}
